% ============================================================================
%  Declaración de uso de IA del Subdocumento 4 (Comunicado 10, sección 7.2).
%  \usoIA{sección}{herramienta}{finalidad}{nivel en texto}{nivel en diagramas}{revisión humana}
%  La revisión humana la firma quien revisó, con nombre y cargo en audIT.
% ============================================================================
\usoIA{Artículo 46}{Claude Opus 5.5 en Claude Code}{Redacción de las respuestas 03 y 53 a 80 y ubicación de la sección que resuelve cada una}{Alto}{Ninguno}{Ignacio Vera, Arquitecto de Solución: correspondencia de cada respuesta con la sección que la resuelve.}
\usoIA{Introducción}{Claude Opus 5.5 en Claude Code}{Redacción inicial del resumen a partir de las decisiones de diseño del equipo}{Alto}{Ninguno}{Ignacio Vera, Arquitecto de Solución: alcance del resumen y su coherencia con las secciones 4.1 a 4.3.}
\usoIA{4 y 4.1.1 a 4.1.26}{OpenAI Codex (GPT-6)}{Contraste con las bases y con los Subdocumentos 3, 5, 6, 7 y 9, corrección de contratos, decisiones y capacidad, y figuras de capas, secuencias y documento sin cobertura}{Alto}{Alto}{Marcel Gutierrez, Arquitecto de Solución: decisiones de arquitectura, contratos, presupuesto de la asignación y supuestos de capacidad de la arquitectura lógica.}
\usoIA{4.1}{Claude Opus 5.5 en Claude Code}{Integración al formato con verificación de citas y versiones, ajuste al formato de las cuatro secuencias y del documento sin cobertura, y diagramas de contextos, integraciones, resiliencia, capa anticorrupción, analítica y modelo táctico generados con un guion sobre los modelos del equipo}{Alto}{Medio}{Ignacio Vera, Arquitecto de Solución: citas a las bases, versiones de cada fabricante y correspondencia de las figuras con el texto.}
\usoIA{4.2}{Claude Opus 5.5 en Claude Code}{Redacción inicial del texto desde las decisiones de diseño del equipo, y diagrama general generado con un guion desde el modelo definido por el equipo}{Alto}{Medio}{Ignacio Vera, Arquitecto de Solución: emplazamiento de cada componente y correspondencia del diagrama general con el Formulario T-11.}
\usoIA{4.2.1}{Claude Opus 5.5 en Claude Code}{Búsqueda de fichas de fabricante, cantidades desde el Caso y redacción inicial}{Alto}{Ninguno}{Ignacio Vera, Arquitecto de Solución: fichas de fabricante y cantidades frente a las poblaciones del Caso.}
\usoIA{4.2.2}{Claude Opus 5.5 en Claude Code}{Fichas del equipo a bordo y cálculo del almacenamiento, y diagrama generado con un guion desde el modelo definido por el equipo}{Alto}{Medio}{Ignacio Vera, Arquitecto de Solución: cálculo del almacenamiento a bordo, especificaciones del equipo y figura.}
\usoIA{4.2.3}{Claude Opus 5.5 en Claude Code}{Conteo del emplazamiento y matriz de trazabilidad}{Alto}{Ninguno}{Ignacio Vera, Arquitecto de Solución: conteo de componentes y matriz de trazabilidad.}
\usoIA{4.2.4}{Claude Opus 5.5 en Claude Code}{Verificación de servicios y niveles en la documentación de Microsoft}{Alto}{Ninguno}{Ignacio Vera, Arquitecto de Solución: servicios, niveles y disponibilidad regional en la documentación de Microsoft.}
\usoIA{4.2.5}{Claude Opus 5.5 en Claude Code}{Plan de direcciones y diagrama de la región primaria generado con un guion desde el modelo definido por el equipo}{Alto}{Medio}{Ignacio Vera, Arquitecto de Solución: plan de direcciones, segmentación y figura de la región primaria.}
\usoIA{4.2.6}{Claude Opus 5.5 en Claude Code}{Redacción inicial de ambientes, promoción y respaldos, y los dos diagramas generados con un guion desde el modelo definido por el equipo}{Alto}{Medio}{Ignacio Vera, Arquitecto de Solución: ambientes, promoción y respaldos, y correspondencia de los dos diagramas.}
\usoIA{4.2.7}{Claude Opus 5.5 en Claude Code}{Enlaces, puntos únicos de falla y tiempos de conmutación con fuente}{Alto}{Ninguno}{Ignacio Vera, Arquitecto de Solución: enlaces, puntos únicos de falla y tiempos de conmutación con su fuente.}
\usoIA{4.2.8}{Claude Opus 5.5 en Claude Code}{Cálculo del dimensionamiento con un guion y diagrama de la reconexión generado con un guion desde el modelo definido por el equipo}{Alto}{Medio}{Ignacio Vera, Arquitecto de Solución: cálculos del dimensionamiento y de la reconexión, y su tratamiento como estimación.}
\usoIA{4.2.9}{Claude Opus 5.5 en Claude Code}{Actualización de los supuestos y brechas del Informe 1}{Medio}{Ninguno}{Ignacio Vera, Arquitecto de Solución: supuestos y brechas frente a las bases.}
\usoIA{4.3}{Claude Opus 5.5 en Claude Code}{Redacción inicial de tipologías}{Alto}{Ninguno}{Ignacio Vera, Arquitecto de Solución: tipología de los centros de datos frente al numeral 6.1 transversal.}
\usoIA{4.3.1}{Claude Opus 5.5 en Claude Code}{Fichas y balance eléctrico con un guion, y plano de la sala y gabinete generados con un guion desde el modelo definido por el equipo}{Alto}{Medio}{Ignacio Vera, Arquitecto de Solución: fichas, balance eléctrico, plano de la sala y gabinete.}
\usoIA{4.3.2}{Claude Opus 5.5 en Claude Code}{Réplicas con fuente, protocolo de conmutación y compromiso de latencia, y diagramas de la recuperación y de la región secundaria generados con un guion desde el modelo definido por el equipo}{Alto}{Medio}{Ignacio Vera, Arquitecto de Solución: réplicas, conmutación, compromiso de latencia y réplica de datos personales en San Bernardo.}
\usoIA{Formulario T-11}{Claude Opus 5.5 en Claude Code}{Armado de las filas desde las fichas y las decisiones de diseño del equipo}{Alto}{Ninguno}{Ignacio Vera, Arquitecto de Solución: cantidades, ubicación y justificación de cada componente.}
